% Discussion, Xingjian Liu. DRAFT, 2026-11-12: written from the forcing probes of MiMo-SFT, MiMo-RL and
% Qwen-Instruct, the core probes, dose and Unseen placement of Qwen-Thinking (keyword count of its reasoning:
% analysis/think_cite.py; placement tests: analysis/place_tests.py), and the RQ1 runs of MiMo-RL and Qwen-Instruct. \tbd marks what waits on Qwen-Thinking,
% MiMo-SFT's RQ1 run and the teammates' sections (RQ1 accuracy, budget sweep, error taxonomy).
% Numbers are Unseen Adversarial unless stated; they come from results/<model>/forcing_tables.md and
% results/<model>/box_swap.md on the data branch rq2/forcing-results.
% Needs xcolor (for \tbd) and the proposal's references.bib.
\providecommand{\tbd}{\textcolor{red}{TBD}}

\section{Discussion}
\label{sec:discussion}

\paragraph{The labels that pull B are mostly the ones the model writes just before it.} In three of
the four models, B is drawn toward A by the labels that come before it in the model's own output,
not by where B sits in the prompt. Qwen-Instruct answers without reasoning, one label after
another. Asking all five objects at once raises its B=A over Single-object by 15.9 points when B is
the last box (AAAAB, Holm $p<10^{-11}$), but only by 2.4 points, not significant, when B is the first
(BAAAA). Given labels pull it in both orders (Section~\ref{sec:forcing}) because it continues from
them, so in both orders A labels end up directly before B. The MiMo models instead rewrite all five
labels in box order, and given labels pull them only in AAAAB, where that rewrite puts four A labels
before B. Inertia in this task is therefore tied to the order of generation, which fits the lower
B=A of Single-object querying, where each box is answered with nothing written before it.
Qwen-Thinking is the exception. Like the MiMo models it rewrites all five labels in box order, so in
BAAAA B is the first label of its answer, yet four given labels still raise its B=A there by 5.7
points (1.5 points per label, Holm $p<0.001$): half its pull on AAAAB, and about as much as
Qwen-Instruct's. With no label before B in its answer, the pull has to come through the prompt or the
reasoning. The reasoning shows one route. Qwen-Thinking often goes back to the given labels as ``the
example in the problem statement'' and decides that all five boxes share one class. By a keyword
count it refers to them in 81\% of the 21 BAAAA images where the labels turn B into A, against 52\%
of those where B stays correct (images where neither call was cut). Answer order is therefore not the
whole story: a reasoning model can carry the given labels into its decision on B even when B comes
first.

\paragraph{In a reasoning model, the earlier labels include the reasoning.} MiMo-RL breaks the
order rule in one place. Under Default its B=A also rises in BAAAA (by 8.7 points over Single-object,
Holm $p=0.001$; Seen: $+1.8$, not significant), although B is the first label of its answer. Its
reasoning, however, goes over all five boxes before the answer is written. Part of this rise is box
confusion. In 12 of the 43 images where B turns into A under Default, the answer names B's class but
writes it on one of the A boxes. We read these 12 reasoning traces: all of them name B's class, and
in 9 the model numbers the boxes in the order it scans them (often from the left) instead of by the
tags drawn in the image. Without these images the rise is 5.1 points. Box confusion does not
explain the forcing effects, though: B's class is written on an A box in 9\% of MiMo-RL's B=A
answers under one-shot teacher forcing on AAAAB, against 12\% under Default. The placement probe also
points at the reasoning. On AAAAB, true labels written at the start of MiMo-RL's answer, after its
reasoning, leave B=A at the Default level ($+3.3$ points, not significant), while the same labels
placed before the reasoning raise it by 13.2. For MiMo-RL, then, the place to look for
self-induced inertia is the reasoning itself. Whether the reasoning names B correctly before the
answer turns to A is the think--answer analysis of Section~\tbd. Qwen-Thinking shows both routes.
Its pull in BAAAA, where B is the first label of its answer, cannot come from labels earlier in
the answer (previous paragraph). Unlike MiMo-RL, it also follows its own answer line: true labels written at
the start of its answer raise its B=A on AAAAB as much as labels in the user turn (44.9\% against
41.6\%, 15.0 points above the one-shot prompt with no labels), and it never re-checks the forced
line. A reasoning model can thus carry this inertia in its reasoning, in its answer, or in both.

\paragraph{RL adds a re-check, but the pull remains.} MiMo-SFT and MiMo-RL share their base model,
so their differences come from RL training. Two things differ clearly. First, only MiMo-RL
re-checks a wrong answer line: after four forced C labels it keeps writing in 62--76\% of items,
typically ``Wait, no, that's wrong'' and a re-check, and by the end of the answer B is C in 8.4\% and
26.0\% of items (AAAAB, BAAAA) instead of 36.8\% and 56.3\% at the first label. MiMo-SFT writes one label and stops (B labelled C in 79.0\% and
86.5\%). This is far more self-correction than the under 10\% that \citet{zhu2026textualinertia}
report for errors injected into the reasoning, probably because here the error contradicts reasoning
the model has just finished. Second, the dose effect of true labels is about half as large after RL:
2.6 points of B=A per label on AAAAB (95\% CI 1.7--3.6) against 5.0 (3.9--6.1) for MiMo-SFT, with
non-overlapping intervals and the same pattern on Seen (2.1 against 5.1). The pull is there before
RL, and RL makes the model less sensitive to the labels before B, not more. It does not remove the
pull: four correct labels still raise MiMo-RL's B=A on AAAAB by 13.2 points over Default. Whether
RL also changes inertia under Default needs MiMo-SFT's RQ1 run (\tbd), and RQ1 (Section~\tbd) gives
the overall accuracy of the two versions.

\paragraph{Rejecting a wrong label is not resisting inertia.} The proposal expected that methods
which detect and repair errors would not help, because in Adversarial items the pull comes from
correct labels. The counterfactual probe supports this. Both MiMo models reject a random wrong class
given in the prompt: B is labelled C in at most 2.4\% of items, and C survives on a given box in
under 5\% of answers. Instead they relabel those boxes themselves, mostly as A (64--74\% of answers),
and B still drifts toward A (MiMo-RL on AAAAB: $+6.3$ points over Default). A wrong class close to A
is accepted more often, but B=A does not change with it. Qwen-Instruct, without reasoning, follows
the wrong class instead (B labelled C in 41.0\% of AAAAB items). Whether a model trusts the labels
it is given and whether its own labels pull it are therefore separate properties, and the MiMo models
control only the first. Reasoning VLMs follow misleading text cues placed beside the
image~\citep{sanchezvillegas2026reasoning}; here the MiMo models reject the cue, and the drift comes
from labels they wrote themselves.

\paragraph{The pull is a consistency prior.} Following the labels is not a failure in itself. On
Homogeneous items, where all five objects share a class, the same four labels raise accuracy on the
fifth box by 6.5 points for MiMo-RL and 10.6 for Qwen-Instruct. The models treat neighbouring boxes
as evidence about each other, which is often right in natural scenes and wrong exactly when one
object differs. Our results suggest two kinds of fix, neither of which we tested beyond the
Single-object baseline. At inference, asking each box separately removes the earlier labels, at five
times the calls. In training, RQ4's per-object reward with an inertia penalty (Section~\tbd) targets
the case our probes isolate, where the earlier labels are correct and B is not.
